\documentclass[11pt]{article}

\usepackage[margin=1in]{geometry}
\usepackage{amsmath,amssymb,amsthm}
\usepackage[T1]{fontenc}
\usepackage{lmodern}
\usepackage{microtype}
\usepackage{xcolor}
\usepackage{hyperref}

\newcommand{\PVol}{\operatorname{PVol}}
\newcommand{\cost}{\operatorname{cost}}
\newcommand{\Path}{\operatorname{\mathbf{Path}}}
\newcommand{\fiber}{\operatorname{fib}}

\newtheorem{observation}{Observation}
\newtheorem{recommendation}{Recommendation}

\title{Notation Note: State-Address Capacity Versus Path-Address Capacity}
\author{}
\date{May 2026}

\begin{document}

\maketitle

\section*{Purpose}

This note records a notational issue in the draft of \emph{Problem agnostic hierarchies, path-volume, and log speed-up in RL}.  The issue is not that the later path-volume theorem is obviously false.  The issue is that an earlier theorem-like preview uses the same symbols \(b_i\) for a state-address capacity, while the later path-volume theorem uses \(b_i\) for an effective path-address capacity.

Those are different mathematical objects.  If they are not distinguished, the paper appears to conflate address size over discovered states with path-space volume.

\section*{The Problem}

The early statement uses positive integers
\[
    b_m,b_{m-1},\ldots,b_0
\]
with an address-capacity condition of the form
\[
    \prod_{i=0}^{m} b_i
    \;\ge\;
    \#\{\text{discovered states}\}
\]
or, in another version,
\[
    \prod_{i=0}^{m} b_i
    \;\ge\;
    |G_t^0|.
\]
This is a statement about capacity to address the discovered base state/action graph.

However, the theorem then describes flat search over a finite path set \(\Omega\) and uses
\[
    |\Omega|
    =
    \PVol_\Omega
\]
as the path-volume term.  In the balanced case it writes
\[
    \prod_i b_i \asymp |\Omega|.
\]
That is a different claim.  In general,
\[
    \#\{\text{discovered states}\}
    \quad\text{and}\quad
    |\Omega|
\]
need not be comparable.  For finite horizon \(K\), the path set can grow exponentially in \(K\) even when the discovered state set is fixed or modest.

\begin{observation}
State-address capacity does not automatically imply path-address capacity.
\end{observation}

The correct bridge is that a tower of state/action partitions can \emph{induce} nested partitions of a finite path set \(\Omega\).  Once that induced path-address system is made explicit, one can state a path-volume theorem.  But the theorem must use the branching factors of the induced path-address system, not merely the branching factors of the state/cell address system.

\section*{Two Different Capacities}

There should be two notationally distinct objects.

\subsection*{State-address capacity}

Let
\[
    b_i^{\mathcal S}
\]
denote a tier-\(i\) state/cell address factor.  This is the kind of quantity that can support a statement such as
\[
    \prod_i b_i^{\mathcal S}
    \;\gtrsim\;
    |\mathcal S_t^0|
\]
or some closely related discovered-graph capacity statement.

This is useful for claims about addressing, navigating, or organizing discovered states and quotient cells.  It is not, by itself, a path-volume statement.

\subsection*{Path-address capacity}

Let
\[
    \Omega\subseteq \Path_{\le K}(G_t^0)
\]
be the finite path set relevant to the learning problem.  Each tier of the tower induces a path-address map
\[
    \pi_{i,*}:\Omega\longrightarrow \Omega_i,
\]
where
\[
    \Omega_i:=\pi_{i,*}(\Omega).
\]
There are coarsening maps
\[
    q_i:\Omega_{i-1}\longrightarrow \Omega_i
\]
between adjacent path-address levels.

A worst-case refinement branching factor can then be defined by
\[
    b_i^\Omega
    :=
    \max_{\eta_i\in\Omega_i}
    \#\left\{
        \eta_{i-1}\in\Omega_{i-1}
        \;:\;
        q_i(\eta_{i-1})=\eta_i
    \right\}.
\]
Equivalently, \(b_i^\Omega\) measures the largest number of tier-\((i-1)\) path-address refinements lying over one tier-\(i\) path-address.

With this notation,
\[
    |\Omega_0|
    \le
    |\Omega_N|\prod_{i=1}^{N}b_i^\Omega.
\]
Since \(\Omega_0=\Omega\), this gives
\[
    |\Omega|
    \le
    |\Omega_N|\prod_{i=1}^{N}b_i^\Omega.
\]
If the coarsest relevant address set \(\Omega_N\) has bounded size, or if its size is absorbed into a top-level factor \(b_N^\Omega\), then the path-address capacity condition becomes
\[
    \prod_{i=0}^{N}b_i^\Omega
    \;\gtrsim\;
    |\Omega|.
\]

\section*{Balanced Addressability}

The phrase ``balanced'' should also refer to the relevant object.

For state addressing, balanced means roughly that the state/cell address capacity is not mostly wasted:
\[
    \prod_i b_i^{\mathcal S}
    \asymp
    |\mathcal S_t^0|.
\]

For path-volume, balanced means that the path-address capacity is not mostly wasted:
\[
    |\Omega_N|\prod_{i=1}^{N} b_i^\Omega
    \asymp
    |\Omega|.
\]
In words: the induced path-address cells and lift fibers should not contain asymptotically huge unused slack relative to the relevant path set.  A distributional or policy-effective version could replace the worst-case \(b_i^\Omega\) by effective support size, entropy, or reward-relevant fiber size under the current rollout distribution.

\section*{Effect On The Later Theorem}

The later path-volume theorem is much closer to the correct statement.  It already phrases the key assumption over the finite path set:
\[
    \prod_i b_i \ge |\Omega|.
\]
It also explains that the numbers \(b_i\) are not graph degrees, but effective address branching factors for the path/refinement problem.

Thus the later theorem is best read as a conditional theorem about path-address capacity.  It is not invalidated by the state/path distinction.  The theorem is assumption-heavy, but coherent:

\[
    \text{if the tower induces a balanced searchable address system on }\Omega,
\]
then
\[
    \cost_{\mathcal T}(\Omega)
    \le
    C\sum_i \log b_i^\Omega
    +
    \sum_i\epsilon_i.
\]
If additionally
\[
    \prod_i b_i^\Omega\asymp |\Omega|,
\]
then
\[
    \sum_i\log b_i^\Omega
    =
    \log\left(\prod_i b_i^\Omega\right)
    =
    \Theta(\log|\Omega|)
    =
    \Theta(\log \PVol_\Omega).
\]

\section*{Recommended Fix}

\begin{recommendation}
Use separate notation for state-address factors and path-address factors.
\end{recommendation}

A clean convention would be:
\[
    b_i^{\mathcal S}
    \quad\text{for state/cell address factors,}
\]
and
\[
    b_i^\Omega
    \quad\text{for path/refinement address factors induced on }\Omega.
\]

Then the early preview theorem should either:

\begin{enumerate}
    \item become a state-address statement only, with no conclusion involving \(\PVol_\Omega\), or
    \item be rewritten to use \(b_i^\Omega\) and the induced path-address maps \(\pi_{i,*}\).
\end{enumerate}

The later path-volume theorem can remain structurally intact if its \(b_i\) are renamed \(b_i^\Omega\), or if the theorem explicitly states that, within that theorem only, \(b_i\) denotes effective path-refinement branching rather than discovered-state address branching.

\section*{Short Verdict}

There is a real notational bug:
\[
    b_i
    \quad\text{is currently doing two jobs.}
\]

There is not necessarily a fatal mathematical bug in the later theorem.  The later theorem is coherent as a conditional path-address theorem.  The early theorem-like preview is the part that currently conflates state-address capacity with path-volume capacity.

\end{document}
